\section{Evaluation}
\label{sec:evaluation}

We evaluate failure-induced workflow delay, the incremental value of Peer HBM
over Peer DRAM, and the conditions under which workflow Slack changes
protection. Discrete-event simulations measure modeled completion time;
separate native-selector observations test selection without assigning it a
serving-time benefit. Negative and null results constrain the design's claims.

\subsection{Methodology}
\label{sec:eval-method}

\paragraph{Simulator and timing model.}
MoonKVSim models dependency release, queues, asynchronous prediction,
checkpoint publication, and recovery on a shared event clock. Transfers contend
on per-rank H2D/D2H links, per-host DRAM and NIC resources, and the fabric.
Pending copies consume capacity; all tensor-parallel (TP) shards must commit
before use. Worker failures remove a TP group; host failures also remove local
DRAM. Fixed token-sequence replays check output counts, epoch fencing, capacity
conservation, and final release, without validating generated token identities.

We use pinned AIConfigurator HYBRID operator estimates\cite{eval2026aiconfigurator}
for dense Qwen3-8B TP1 and Qwen3-32B TP4 on A100-SXM configurations.
Each 196-shape grid spans 49 prefixes (0--98,304, step 2,048), prefill extensions
of 1/64/256 tokens, and one-token decode. Bounded interpolation supplies forward
durations; decode steps sum individually. Empirical, mixed, and silicon-estimate
provenance is retained. This is not end-to-end calibration. Foreground execution
is FCFS with one active invocation per data-parallel worker, so results do not
establish production throughput or continuous-batching behavior.

\paragraph{Workload and held-out split.}
The primary workload is a constructed, evidence-provided FanOutQA
replay\cite{eval2024fanoutqa} with explicit dependencies. The retained source
contains 160 workflows from 310 development questions, selected using a
previously frozen task-ID order and context guard. Keeping complete workflows
within the 40,960-token context bound leaves 74 workflows. A predeclared
workflow-ID hash split gives 12 training workflows (95 tasks) and 62 evaluation
workflows (503 tasks). Each original DAG executes once; each task is a distinct,
single-invocation Agent. Prompts use pinned Qwen tokenizers; outputs are
tokenized gold answers. A declared 512-token output limit determines reservation;
eventual output length and future invocations are unavailable to the policy.

FanOutQA does not validate naturally occurring multi-call histories. Separate
constructed four-Agent fork--join interventions cover skewed/balanced branches,
same-Agent history reuse, long decode with changing synthetic demand, and
shared hosts. Seeds 7/19/41 vary token identities and prescribed length jitter.
Each main case has eight workflows and four TP1 workers; shared-host cases
place two workers on each of two hosts. The declared output cap is 16,384.
These interventions are not additional natural-trace samples.

\paragraph{Resources and predictions.}
The primary deployment has eight workers on independent hosts: eight GPUs for
8B and 32 for 32B. Each GPU has 80\,GiB HBM less profile weight/communication
allocations and a 4\,GiB runtime allowance; each host has 256\,GiB Peer DRAM.
Configured effective H2D/D2H, NIC/fabric, and DRAM bandwidths are 32, 25, and
100\,GB/s. Peer HBM starts at 128\,MiB per rank, grows by 64\,MiB after stable
windows, and is capped at 4\,GiB, with a 2\,GiB foreground reserve.
Blocks contain 64 tokens; checkpoints and prefill chunks span 256 tokens.
We set $K=16$, $\beta=0.9$, $\epsilon=0.1$, and $\lambda=10$\,ms.
Metadata uses 16 selector-state plus eight counter bytes per block/layer/rank.

The predictor combines submitted-prompt work with training role-median output
length. Unstarted Agents use training role-duration priors; the guard is the
training 90th-percentile absolute residual. This estimates invocation service,
excluding queues, gaps, and future calls. Subtracting age from submission can
therefore invalidate the estimate and force workflow-wide zero Slack.
Reconstruction estimates and execution share a profile, an idealization rather
than independent prediction validation. The offline depth rule uses the
checkpoint-aligned training-median context (8,704 tokens), full dense demand
including metadata, and uniform per-layer compute. It fixes $L=11/1$ for
8B/32B before evaluation. Constructed controls use declared priors and $L=4$.

\paragraph{Comparisons and failures.}
Five policies share one runtime: restart without checkpointing; \emph{All-DRAM},
which removes the exposure threshold while retaining exposure ordering;
\emph{DRAM}, with \system's workflow selection but no extra HBM;
\emph{Static-HBM}, which starts at 128\,MiB, disables quota growth, and still
allows foreground enforcement to shrink it; and full \system. These are
controlled variants, not complete prior-system reproductions. Comparing with
DRAM isolates HBM from avoided reconstruction.

Arrivals divide mean training workflow work by eight workers and nominal load
$\rho\in\{0.3,0.65,0.9\}$, a generation parameter rather than measured
utilization. Each model/load/policy has one healthy run and five failures:
180 runs. Five equally spaced source-order workflows alternate sink and middle
targets, interrupted at 80\% of the same invocation's total logical tokens.
Short outputs place four faults in prefill and one in decode. The separate
135-run constructed matrix instead targets 80\% of prefill and decode
\emph{phase} progress, which are different positions. Detection is 1\,ms and
failed capacity remains lost. Primary work-rate proxies after losing one of
eight workers are $8\rho/7$, reaching 1.03 at the highest load; these are not
measured utilization. Logical alignment fixes target identity and
token position, but can change wall time or collateral victims. Splitting a
modeled prefill at a fault can change its profiled compute cost; paired
penalties include this effect as well as recovery and scheduling.

\paragraph{Metrics and reproducibility.}
WCT is last-Agent completion minus external workflow arrival. Co-primary
endpoints are the target workflow's
$\Delta\mathrm{WCT}=\mathrm{WCT}_{F1}-\mathrm{WCT}_{F0}$ and the interrupted
Agent's failure-to-first-new-token latency, preserving fixed identities and
sink/non-sink strata. The latter decomposes into detection, retry queue, and
service from admission through the first new output token. The all-workflow
mean separately captures collateral scheduling and permanent capacity loss.
Artifacts retain absolute healthy/failure WCT, checkpoint lag, byte and memory
accounting, co-victims, and per-identity recovery; selected runs retain full
transfer intervals. Five fixed faults are descriptive scenarios; workflows
within a run are dependent. Reusing F0 does not create new healthy samples.
The retained 62-workflow p50/p95/p99 are interpolated descriptions, not reliable
population tails. Deterministic reruns receive no confidence intervals.

All 483 current simulations use the same runtime version, sealed per campaign:
180 primary, 135 structural,
126 sensitivity, 32 capacity/Slack, and ten post-hoc mechanism runs.
It selects over causally visible KV, counts locally recomputed tail KV, restores
committed histories for queued continuations, and cancels unused restoration at
accepted EOS. The 74 runtime regressions and per-campaign audits accompany
matched-input replacements of earlier versions; those historical runs are
retained separately and are not additional samples.

\subsection{Primary Workflow Replay}
\label{sec:eval-primary}
Figure~\ref{fig:eval-primary} retains every target workflow, separating sink
and non-sink identities. For 8B, DRAM reduces mean target penalty from
0.530 to 0.062\,s at $\rho=0.3/0.65$, but increases it from 1.994 to
2.070\,s at $\rho=0.9$. The cohort means move differently: at $\rho=0.65$
DRAM worsens 0.744 to 0.766\,s, whereas at $\rho=0.9$ it improves 3.956 to
3.587\,s. For 32B at $\rho=0.3$, the mean target penalty improves from
1.427 to 1.124\,s, but the worst identity worsens from 3.876 to 5.299\,s.
Its retry changes workers: queueing adds 1.453\,s while service saves 29\,ms.
Placement and reconstruction effects therefore both contribute to WCT.

Table~\ref{tab:agent-failure-latency} reports the interrupted Agent's separate
endpoint. Retry queue accounts for 62.22--90.63\% of DRAM's mean latency
across the six model/load cells; retries cannot preempt a busy Peer.
The artifact retains
each identity's queue/service decomposition and sink/non-sink summaries.
Three sink targets fail at only 257--356 logical tokens with a 256-token
checkpoint, limiting the reconstruction work they can avoid.

The target-workflow HBM increment ranges from $-0.0026$ to 1.442\,ms for 8B
and from 0 to 0.137\,ms for 32B across the 15 model-specific fault conditions.
Static-HBM and \system produce identical WCT for every workflow in this
matrix despite different allocation and transfer behavior; these runs do not
demonstrate a WCT advantage from quota adaptation. Healthy WCT is also
identical across all five policies: its mean ranges from 6.23 to 16.65\,s
across model/load conditions. Backup traffic is nonzero; the unchanged WCT
reflects this model's execution and interference assumptions.

\begin{table}[t]
\centering
\setlength{\tabcolsep}{3pt}
\renewcommand{\arraystretch}{1.08}
\caption{Mean failure-to-first-new-token latency over five fixed identities.
HBM gain is DRAM minus AegisKV; positive is better. No confidence intervals.}
\label{tab:agent-failure-latency}
\begin{tabular}{@{}lrrrr@{}}
\hline
Model & $\rho$ & Restart & DRAM & HBM gain \\
 & & (s) & (s) & (ms) \\
\hline
Qwen3-8B & 0.3 & 0.956 & 0.488 & 0.220 \\
Qwen3-8B & 0.65 & 0.956 & 0.488 & 0.219 \\
Qwen3-8B & 0.9 & 2.694 & 1.970 & 0.508 \\
Qwen3-32B & 0.3 & 1.910 & 1.607 & 0.027 \\
Qwen3-32B & 0.65 & 1.144 & 0.551 & 0.027 \\
Qwen3-32B & 0.9 & 2.222 & 1.253 & 0.036 \\
\hline
\end{tabular}
\end{table}

\begin{figure*}[t]
  \centering
  \includegraphics[width=\textwidth]{figures/eval-primary-wct.pdf}
  \caption{Target-workflow penalty for five fixed identities at each model/load.
  S/N denote sink/non-sink Agents. Each marker subtracts own-policy healthy WCT
  for the same workflow; offsets separate coincident policies. Logical targets
  match across policies; negative penalties remain visible. No confidence
  intervals. Nominal load is an arrival-generation parameter.}
  \label{fig:eval-primary}
  \Description{Six panels compare target-workflow paired delay at five fixed sink or non-sink failure identities for two models and three nominal loads. Five policies use offset markers; negative penalties are retained.}
\end{figure*}

DRAM versus All-DRAM isolates the exposure threshold. Healthy transferred
backup bytes fall by 10.21\%, 5.70\%, and $-0.095$\% at increasing 8B loads;
32B savings are 10.81\%, 7.98\%, and 1.16\%, including cancellation bytes.
Cohort WCT means match, although two workflows exchange $\pm4.113$\,s in
one 32B high-load fault. In the two fully logged \system F1 runs at
$\rho=0.65$, fallback covers 3,122/5,586 and 3,605/6,518 unfinished-workflow
control-window observations (55.89\%/55.31\% for 8B/32B), all because prediction
minus submission age is nonpositive. These are neither healthy-run percentages
nor time fractions. Frequent zero-Slack fallback limits what this replay can
establish about workflow-aware prioritization.

A small recovery gain can change subsequent routing. In the 8B high-load
second fault, HBM advances the target workflow by 1.10\,ms. This moves another
completion across a later fork's submission, changes the Router's outstanding
request counts, and assigns its branches to different workers. One other
workflow finishes 3.41\,s later; eight workflows gain 1.10\,ms each, including
the target, and 53 are unchanged. The 62-workflow WCT sum therefore increases
by 3.40\,s. Fault time, failed worker, active co-victims, and queued migration
match. Saved submission, completion, and worker records reproduce the routing
change without treating unrecorded admission delay as measured queue time.
This counterexample limits a direct interpretation of local $G_a$ as a global
workflow benefit.

\subsection{Workflow Structure and Recovery}
\label{sec:eval-structure}
Across the 135 constructed runs, DRAM worsens six of ten scenario/phase
mean target penalties and two cohort means (Figure~\ref{fig:eval-structure}).
Both worsen for skewed prefill (14.43/2.89\,s, target/cohort) and long-decode
decode (21.24/6.63\,s); the latter's seed 19 worsens by 49.70/15.34\,s.
Long-decode prefill also worsens the target by 30.62\,s. Extra HBM changes
the ten mean target penalties by less than 0.875\,ms, and by below one
microsecond in long-decode cases. There is no uniform checkpoint improvement
or substantial incremental HBM benefit.

\begin{figure*}[t]
  \centering
  \includegraphics[width=\textwidth]{figures/eval-structural-wct.pdf}
  \caption{Target-workflow penalty in five constructed four-Agent workloads,
  eight workflows per run. Gray segments connect matched seeds; colored ranges
  show min/max and horizontal marks means. Each F1 subtracts own-policy F0.
  Agent, second invocation and logical position match across policies;
  prefill/decode positions differ. Negative penalties and overlapping
  DRAM/AegisKV points remain visible.}
  \label{fig:eval-structure}
  \Description{Two panels show target-workflow delay for prefill and decode faults across five constructed scenarios, with three matched seed points per policy. Negative penalties remain visible. Shared-host failures have large variation; checkpointing can worsen the target workflow.}
\end{figure*}

Endpoint choice can reverse the conclusion. In shared-host prefill/decode
cases, DRAM worsens the mean target penalty by 30.00/41.01\,s while reducing
the all-workflow mean penalty by 16.11/19.01\,s. History-reuse prefill likewise
worsens the target by 12.05\,s but improves the cohort mean by 3.33\,s.
Permanent capacity loss, queue placement, and completion-dependent routing
redistribute delay across workflows; a local gain $G_a$ is not a global WCT
guarantee. The audit separately records 108 active interrupted recoveries and
279 queued-history restorations, without counting the latter as active
co-victims. All per-seed target/cohort endpoints and recovery decompositions
remain in the artifact.

One matched history-prefill case (seed 19) makes the reversal explicit.
DRAM shortens recovery service from 567.66 to 42.54\,ms with the same 1.060\,s
retry queue. The final left call is submitted 0.525\,s earlier but lasts
26.174\,s longer; the join adds another 17.681\,s of invocation latency.
Target WCT worsens by 43.330\,s while the other seven workflows improve by
81.811\,s in total. Recorded placement and timestamps support this decomposition;
ordinary admission is not recorded, so these latter latency changes are not
reported as measured queue time.

The long-decode control uses two synthetic four-block banks and dense Qwen
compute. Its $K=16$ fits both banks, so it cannot establish robustness to native
hotspot drift under an insufficient budget. Larger 32/128-workflow simulator
CPU/RSS controls are retained under the historical runtime; they are neither
current recovery results nor physical controller-overhead measurements.

\subsection{Operating Range and Controls}
\label{sec:eval-sensitivity}
The 126 one-factor controls use four complete held-out workflows (52 Agents),
fixed arrivals/priors, and the longest-prompt Agent in the third workflow:
20,283 prompt tokens, interrupted at token 16,228. Input length selects this
target before outcomes. Each condition subtracts its own F0; shared center
runs are not additional samples. This is one fixed fault, not a population.

In Figure~\ref{fig:eval-resources}, reducing H2D from 64 to 8\,GB/s raises the
DRAM target penalty from 24.02 to 283.95\,ms; HBM saves another 0.72/5.78\,ms.
At 32\,GB/s, restart/DRAM/\system penalties are 1.8023/0.0610/0.0595\,s,
with a 1.44\,ms HBM increment. At 0.25/1\,GiB Peer DRAM, the target loses its
checkpoint and both backup policies revert to restart; peak host occupancy
reaches 98.45\%. This is a capacity boundary, not unconstrained DRAM operation.
A zero/32\,MiB HBM cap gives the same target penalty as DRAM.

A 32-run control crosses 2/4/8/256\,GiB DRAM with DRAM, All-DRAM, \system,
and \system with zero DAG Slack, affecting both exposure and HBM gain.
Current and zero-Slack variants have identical individual workflow WCT in
all F0/F1 conditions. Current reduces healthy backup bytes by 4.58--4.66\%,
including completed checkpoint/hot placement and bytes served before cancellation.
The target checkpoint needs 2.215\,GiB: it cannot fit alone at 2\,GiB and
remains at 16,128 tokens for every variant at larger budgets. This cohort shows
traffic savings, not a competition-induced WCT benefit.

To exercise protection competition, a post-hoc ten-run diagnostic uses one
constructed fork--join, declared branch priors giving 0/80\,s Slack, and no
Peer-HBM backup. Its 1.25\,GiB shared peer budget fits either branch checkpoint
individually but not both. At the fixed critical-branch fault, DAG Slack retains
its 1,280-token checkpoint; zero Slack instead retains the expensive noncritical
branch's 8,448-token checkpoint. With matching fault time and no retry queue,
workflow penalties are 23.97/134.79\,ms. At 4\,GiB both variants retain the
critical checkpoint and yield 23.97\,ms. This validates a conditional selection
mechanism, not prediction quality or its prevalence in natural workloads.

\begin{figure*}[t]
  \centering
  \includegraphics[width=\textwidth]{figures/eval-resource-sensitivity.pdf}
  \caption{One fixed prefill fault on four complete FanOutQA workflows,
  Qwen3-8B TP1, eight workers. Points are target-workflow F1 minus own-condition
  F0 WCT; no confidence interval is implied. $L=11$, $K=16$, and initial HBM
  quota $\min(128\,\mathrm{MiB},\mathrm{cap})$. Delay, bandwidth and DRAM axes
  are logarithmic; the HBM axis is linear from 0 to 32\,MiB and logarithmic
  thereafter.}
  \label{fig:eval-resources}
  \Description{Three panels vary H2D bandwidth, host DRAM and per-rank HBM. Faster H2D lowers recovery delay. Small DRAM budgets lose the target checkpoint; low HBM quotas match DRAM-only.}
\end{figure*}

Other controls expose limits. Prediction multipliers 0.5/1/2 and
$\lambda=0/0.01/1$\,s leave the target penalty unchanged. With a 100\,ms TTL,
all 52 replies at 200\,ms are stale, yet changing latency from 0.1 to 200\,ms
also leaves that penalty unchanged. These non-crossing controls check code
paths, not prediction accuracy or robustness; detection controls use 1/10/100\,ms.

DMA slowdown coefficients 0.05/0.1 increase healthy mean WCT by 3.33/6.66\%,
from 8.398\,s to 8.678/8.957\,s; even F1 cohort WCT sums worsen by 0.39/0.74\,s.
For a hypothetical mixture of F0 and this exact
F1 with probability $p$, advantage is $pD-H$: healthy cost $H$ and reduction
in paired fault penalty $D$ must use the same objective. At 0.1, break-even
thresholds are 0.141 for the target and 1.494 for the four-workflow sum; thus
no feasible $p\in[0,1]$ offsets cohort cost. These are not measured failure rates.
Finally, $Z=64$ also limits forwards to 64 tokens, whereas $Z=1024$ retains
256-token chunks. The former raises even restart's penalty to 17.71\,s:
a coupled compute/checkpoint intervention. Full controls and adverse outcomes
remain in the artifact.

\subsection{Native Selection and its Limits}
\label{sec:eval-native}
We separately replay sealed NOSA-8B, NSA-8B, and MiniCPM4.1-8B selectors on the
same three PG19 documents. We union blocks across KV groups per query, start
counters at query zero, and apply decay, hysteresis, and the earlier-block tie
rule every 256 tokens. All 13,824 checked sets match the deployed ranking
function. This assumes continuous protection and immediate checkpoint visibility;
it does not reproduce publication gaps or native chunked-prefill execution.

\begin{figure}[!htbp]
  \centering
  \includegraphics[width=\columnwidth]{figures/eval-native-topk.pdf}
  \caption{Native distinct-history coverage over the first four layers and
  next 256 queries at checkpoints 32,000/32,256. Each document's ratio sums
  covered/requested block counts over layers and checkpoints. Points average
  three documents; whiskers show their range. The future-demand oracle uses
  the same per-layer block budget $K$. $\Delta$ is Current minus Recency at
  $K=64$, in percentage points.}
  \label{fig:eval-native-topk}
  \Description{Three native-model panels compare current selection, recency, and a future-demand oracle as the block budget grows. Current beats recency for NOSA and NSA at budget 64 but loses for MiniCPM.}
\end{figure}

For NOSA/NSA/MiniCPM respectively, $K=16,L=4$ gives Current distinct coverage
of 8.31\%, 4.20\%, and 4.06\%. Repeated query--block event coverage is
15.70\%, 11.46\%, and 5.37\%, versus Recency's 15.18\%, 15.44\%, and 11.45\%
and the oracle's 15.83\%, 16.15\%, and 11.45\%. These observations do not
support ``few misses'' at this budget. They also differ from the dense WCT
models and $L=11/1$ configurations.

At $K=64$, Current distinct coverage rises to 32.77\%, 16.76\%, and 15.45\%
(Figure~\ref{fig:eval-native-topk}), improving on Recency by
2.36, 0.97, and $-1.07$ percentage points. The oracle reaches only
33.25\%, 16.79\%, and 16.53\%, showing that the budget itself limits coverage.
Event coverage is 59.90\%, 36.04\%, and 22.35\% for Current;
33.77\%, 34.37\%, and 33.77\% for Recency; and 63.32\%, 53.60\%, and 42.69\%
for the oracle. MiniCPM loses on both measures. The denominators cannot be
added or interpreted as runtime cache-hit rates, and this three-document
diagnostic does not establish a WCT advantage.

\subsection{Scope of the Evidence}
\label{sec:eval-limits}
Unmodeled CPU/kernel costs and batching, and uncalibrated DMA interference,
preclude physical-overhead claims. Immediate continuations exercise queued
history recovery without tool/think idle gaps.

For one Agent's dense checkpoint, HBM holds at most $LK$ historical blocks per
rank; other demanded blocks need H2D. This partial residency does not exercise
Equation~\ref{eq:depth}'s no-wait first-$L$-layers premise. Bytes/bandwidth
explains local service savings, not global WCT under rerouting. FCFS also
does not establish quota adaptation's benefit under changing batch memory pressure.

Asynchronous publication can leave more than $Z$ uncommitted tokens.
A historical, pre-freeze linear-cost control accepts 6,247 tokens with
committed frontier 1,024: a 5,223-token tail despite $Z=256$.
Likewise, protection tolerance and modeled depth are parameters, not
per-workflow delay or transfer-hiding guarantees.
